\chapter{Compactness}

\begin{center}\itshape
A continuous function on a closed interval reaches its maximum; on an open interval it may not. A bounded sequence of
real numbers has a convergent subsequence; a bounded sequence of continuous functions may not. What do the good
cases share, and how far does the property reach?
\end{center}

A drunkard wanders forever along a closed segment of road. Must he come back, again and again, arbitrarily close to some
spot? On a segment, yes: that is the Bolzano--Weierstrass theorem, and the reason Tome~0 could find maxima of continuous
functions. On an infinite road, no; he may simply walk away. And in an infinite-dimensional ``segment'', the closed unit
ball of a space of functions, the answer is again no, for a more subtle reason: there are infinitely many directions to
walk away in. Compactness is the name of the property that keeps the drunkard home.

Finite sets are easy: a function on a finite set has a largest value, finitely many positive numbers have a positive
minimum. Compactness is what lets an infinite space behave like a finite one in such arguments, and its history is a slow
search for the right definition. Heine (1872) and Borel (1895) noticed that any cover of $[a,b]$ by open intervals can be
reduced to a finite one; Fr\'echet coined the word ``compact'' in 1906; Alexandroff and Urysohn (1923) made the cover
property the definition. It looks strange at first, and it is exactly what the proofs need. Unlike the sequential
definition, it survives infinite products.

This chapter first keeps the last promises of Tome~0: extreme values, Heine's uniform continuity, Rolle and the mean value
theorem. Then come the three great compactness theorems of analysis. Tychonoff's theorem says that \emph{any} product of compact
spaces is compact; in 1950 John Kelley showed that it is in fact equivalent to the axiom of choice. Arzel\`a--Ascoli describes the
compact sets of continuous functions, and Stone--Weierstrass explains when a family of functions can approximate every
continuous function.

\section{Compact spaces}

\begin{definition}
An \emph{open cover} of a topological space $X$ is a family of open sets whose union is $X$. The space $X$ is \emph{compact} if every open cover
has a finite subfamily that still covers $X$ (a \emph{finite subcover}). A subset $K\subseteq X$ is compact if it is compact for the subspace
topology, that is, if every cover of $K$ by open sets of $X$ has a finite subcover.
\end{definition}

Passing to complements: $X$ is compact iff every family of closed sets with empty intersection has a finite subfamily with
empty intersection. Say that a family of sets has the \emph{finite intersection property} if every finite subfamily has non-empty
intersection. Then $X$ is compact iff \emph{every family of closed sets with the finite intersection property has non-empty
intersection}. For instance a decreasing sequence $F_1\supseteq F_2\supseteq\cdots$ of non-empty closed sets of a compact space has a common point.

\begin{think}
Is $[0,1[$ compact? Is $\Z$ with the discrete topology? Is any finite space?
\end{think}

The cover of $[0,1[$ by the sets $[0,1-\frac1n[$ has no finite subcover, nor does the cover of $\Z$ by singletons. Every finite space is compact, and
so is any space with finitely many open sets.

\begin{proposition}\label{prop:compactbasic}
\begin{enumerate}[label=(\alph*)]
\item A closed subset of a compact space is compact.
\item A compact subset of a Hausdorff space is closed.
\item A continuous image of a compact space is compact.
\item A continuous bijection from a compact space onto a Hausdorff space is a homeomorphism.
\item A compact Hausdorff space is normal: two disjoint closed sets have disjoint open neighbourhoods.
\end{enumerate}
\end{proposition}
\begin{proof}
(a) If $F\subseteq X$ is closed and $\mathcal U$ covers $F$ by open sets, then $\mathcal U\cup\{X\setminus F\}$ covers $X$; a finite subcover, with $X\setminus F$ removed,
covers $F$.

(b) Let $K$ be compact in a Hausdorff space $X$, and $x\notin K$. For each $y\in K$ choose disjoint open $U_y\ni x$ and $V_y\ni y$. Finitely many $V_{y_1},\dots,V_{y_m}$
cover $K$, and $U=U_{y_1}\cap\cdots\cap U_{y_m}$ is an open neighbourhood of $x$ that misses all of them, hence misses $K$. So $X\setminus K$ is open.

(c) If $f:X\to Y$ is continuous and $\mathcal V$ covers $f(X)$ by open sets, the sets $f^{-1}(V)$, $V\in\mathcal V$, cover $X$; finitely many suffice, and the
corresponding $V$'s cover $f(X)$.

(d) Let $f:X\to Y$ be the bijection. For $F$ closed in $X$, $(f^{-1})^{-1}(F)=f(F)$ is compact by (a) and (c), hence closed by (b). So $f^{-1}$ is continuous.

(e) First let $A$ be a point $a$ and $B$ closed, $a\notin B$. The proof of (b), applied to the compact set $B$ (by (a)), gives open sets $U\ni a$ and
$V=V_{y_1}\cup\cdots\cup V_{y_m}\supseteq B$ that are disjoint. Now let $A$ be closed, disjoint from $B$. For each $a\in A$ take such disjoint $U_a\ni a$ and $V_a\supseteq B$.
Finitely many $U_{a_1},\dots,U_{a_r}$ cover the compact set $A$; then $U=\bigcup_iU_{a_i}\supseteq A$ and $V=\bigcap_iV_{a_i}\supseteq B$ are open and disjoint.
\end{proof}

Part~(d) is used constantly: to check that a continuous bijection between compact spaces is a homeomorphism, nothing needs
to be done. This is how $\R/\Z\cong S^1$ in Chapter~3 can be seen without computing inverses: $[0,1]$ maps continuously onto both.
Part~(e) and the remark after Proposition~\ref{prop:urysohn} give Urysohn functions on compact Hausdorff spaces.

\begin{proposition}\label{prop:finiteproduct}
If $X$ and $Y$ are compact, so is $X\times Y$. Hence a finite product of compact spaces is compact.
\end{proposition}
\begin{proof}
It is enough to treat covers by \emph{basic} open sets $U\times V$: given any open cover $\mathcal W$, the basic sets contained in some member of $\mathcal W$
cover $X\times Y$, and if finitely many of them cover, so do the finitely many members of $\mathcal W$ that contain them.

So let $\mathcal B$ be a cover of $X\times Y$ by basic open sets. Fix $x\in X$. The slice $\{x\}\times Y$ is homeomorphic to $Y$, hence compact, and finitely many
$U_1\times V_1,\dots,U_m\times V_m$ in $\mathcal B$ cover it; we may assume each meets the slice, so $x\in U_j$. Then $U_x=U_1\cap\cdots\cap U_m$ is an open neighbourhood of
$x$, and the whole \emph{tube} $U_x\times Y$ is covered by the same $m$ sets: a point $(x',y)$ of the tube has $(x,y)\in U_j\times V_j$ for some $j$, and $x'\in U_x\subseteq U_j$.
Finitely many $U_{x_1},\dots,U_{x_r}$ cover $X$, and the corresponding finitely many members of $\mathcal B$ cover $X\times Y$.
\end{proof}

Now the payoff. A closed segment is compact; a closed box in $\R^n$ is a product of segments; and the drunkard is
trapped.

\begin{theorem}[Heine--Borel]\label{thm:heineborel}
A closed interval $[a,b]$ of $\R$ is compact. A subset of $\R^n$ is compact iff it is closed and bounded.
\end{theorem}
\begin{proof}
Let $\mathcal U$ be an open cover of $[a,b]$ and $S$ the set of $x\in[a,b]$ such that $[a,x]$ is covered by finitely many members of $\mathcal U$. Then $a\in S$; let
$s=\sup S$ (least upper bound property, Theorem~\ref{thm:reals}). Some $U\in\mathcal U$ contains $s$, and with it $]s-\delta,s+\delta[\cap[a,b]$ for some $\delta>0$. There is
$x\in S$ with $x>s-\delta$; adding $U$ to the finitely many sets covering $[a,x]$ covers $[a,\min(s+\frac\delta2,b)]$. So $\min(s+\frac\delta2,b)\in S$, which forces
$s=b$ and $b\in S$.

Let $K\subseteq\R^n$ be closed and bounded, say $K\subseteq[-M,M]^n$. The topology of $\R^n$ is the product topology (Proposition~\ref{prop:product}(b)), so
$[-M,M]^n$ is compact by Proposition~\ref{prop:finiteproduct}, and $K$, closed in it, is compact. Conversely a compact $K$ is closed by
Proposition~\ref{prop:compactbasic}(b), and bounded because the open balls $B(0,k)$, $k\ge1$, cover it and finitely many suffice.
\end{proof}

\begin{corollary}[extreme values]\label{cor:extreme}
A continuous function $f:K\to\R$ on a non-empty compact space is bounded and attains its maximum and its minimum.
\end{corollary}
\begin{proof}
$f(K)$ is compact in $\R$ (Proposition~\ref{prop:compactbasic}(c)), so closed and bounded. Its supremum is a limit of points of $f(K)$, hence belongs to
$f(K)$; the same for the infimum.
\end{proof}

Tome~0 admitted this for $K=[a,b]$; it is now proved. Here is a first use, promised in Chapter~2.

\begin{corollary}\label{cor:norms}
All norms on $\R^n$ (or $\C^n$) are Lipschitz equivalent. Hence every finite-dimensional normed space is complete, and its compact subsets are its closed
bounded subsets.
\end{corollary}
\begin{proof}
Let $N$ be a norm on $\R^n$ and $(e_i)$ the standard basis. Then $N(x)\le\sum_i|x_i|N(e_i)\le C\|x\|_\infty$ with $C=\sum_iN(e_i)$, so $|N(x)-N(y)|\le N(x-y)\le C\|x-y\|_\infty$:
$N$ is continuous for $\|\cdot\|_\infty$. The sphere $S=\{x:\|x\|_\infty=1\}$ is closed and bounded, hence compact, and $N>0$ on it, so $N\ge c$ on $S$ for some
$c>0$ (Corollary~\ref{cor:extreme}). By homogeneity $c\|x\|_\infty\le N(x)\le C\|x\|_\infty$ for all $x$. The complex case is the real case in $\R^{2n}$. A finite-dimensional
normed space is isomorphic to some $\R^n$ or $\C^n$ by the choice of a basis, which transports the norm to a norm $N$ on $\R^n$ or $\C^n$; completeness
(Proposition~\ref{prop:complete}(b)) and the description of compact sets (Theorem~\ref{thm:heineborel}) are unchanged by Lipschitz equivalence.
\end{proof}

\section{Compact metric spaces}

In metric spaces, compactness has two other faces, through sequences and through finite nets.

\begin{definition}
A metric space $X$ is \emph{sequentially compact} if every sequence has a convergent subsequence. It is \emph{totally bounded} if for every $\varepsilon>0$ it is
covered by finitely many balls of radius $\varepsilon$; their centres form an \emph{$\varepsilon$-net}.
\end{definition}

\begin{figure}[ht]
\centering
\fig{ch04-net}
\caption{Total boundedness: for every $\varepsilon$, finitely many balls of radius $\varepsilon$ cover the space. Their centres (the dots) form an $\varepsilon$-net. A
compact set is one that can be ``seen to precision $\varepsilon$'' with finitely many points, for every $\varepsilon$, and is complete.}
\label{fig:net}
\end{figure}

\begin{theorem}\label{thm:compactmetric}
For a metric space $X$ the following are equivalent:
\begin{enumerate}[label=(\roman*)]
\item $X$ is compact;
\item $X$ is sequentially compact;
\item $X$ is complete and totally bounded.
\end{enumerate}
Moreover, if $X$ is compact, every open cover $\mathcal U$ has a \emph{Lebesgue number}: some $\delta>0$ such that every ball of radius $\delta$ lies inside a member of $\mathcal U$.
\end{theorem}
\begin{proof}
(i)$\Rightarrow$(ii). Let $(x_n)$ be a sequence. Call $x$ a \emph{cluster point} if every ball around $x$ contains $x_n$ for infinitely many $n$. If there is none, every $x$
has an open ball $B_x$ containing $x_n$ for only finitely many $n$; finitely many $B_x$ cover $X$, so only finitely many indices $n$ exist, which is absurd. So
there is a cluster point $x$, and choosing $n_1<n_2<\cdots$ with $d(x_{n_k},x)<\frac1k$ gives a subsequence converging to $x$.

(ii)$\Rightarrow$(iii). A Cauchy sequence with a convergent subsequence converges (to the same limit), so $X$ is complete. If $X$ were not totally bounded,
for some $\varepsilon>0$ no finite family of $\varepsilon$-balls would cover $X$; choose $x_1$, then $x_2\notin B(x_1,\varepsilon)$, then $x_3\notin B(x_1,\varepsilon)\cup B(x_2,\varepsilon)$, and so on. Then
$d(x_m,x_n)\ge\varepsilon$ for $m\ne n$, and no subsequence is Cauchy.

(iii)$\Rightarrow$(ii). Let $(x_n)$ be a sequence. Finitely many balls of radius $1$ cover $X$, so one of them, $B_1$, contains $x_n$ for infinitely many $n$: let $N_1$ be the set
of these $n$. Finitely many balls of radius $\frac12$ cover $X$, so one of them, $B_2$, contains $x_n$ for infinitely many $n\in N_1$: call this set $N_2$. Continuing gives infinite
sets $N_1\supseteq N_2\supseteq\cdots$ with $x_n\in B_k$ for $n\in N_k$. Choose $n_1<n_2<\cdots$ with $n_k\in N_k$. For $j\ge k$, $x_{n_j}$ and $x_{n_k}$ both lie in $B_k$, so
$d(x_{n_j},x_{n_k})<\frac2k$: the subsequence is Cauchy, and converges since $X$ is complete.

(ii)$\Rightarrow$(i). First the Lebesgue number, assuming (ii). If there is none, for each $n$ some ball $B(x_n,\frac1n)$ lies in no member of $\mathcal U$. A subsequence
$x_{n_k}$ converges to some $x$; $x\in U$ for some $U\in\mathcal U$, and $B(x,r)\subseteq U$ for some $r>0$. For $k$ large, $d(x_{n_k},x)<\frac r2$ and $\frac1{n_k}<\frac r2$, so $B(x_{n_k},\frac1{n_k})\subseteq B(x,r)\subseteq U$,
a contradiction. Now let $\delta$ be a Lebesgue number of $\mathcal U$. Since (ii) implies total boundedness, finitely many balls of radius $\delta$ cover $X$, and each lies
in a member of $\mathcal U$: these finitely many members cover $X$.
\end{proof}

So a compact metric space is bounded, and \emph{separable}: the union of finite $\frac1k$-nets, $k\ge1$, is a countable dense set. In $\R^n$, Theorem~\ref{thm:compactmetric}
gives the Bolzano--Weierstrass theorem: every bounded sequence has a convergent subsequence (its closure is closed and bounded, hence compact).

\begin{think}
Is the closed unit ball of $C([0,1])$, with the sup norm, compact? It is closed, bounded and complete.
\end{think}

It is not. Let $f_n$ be the ``tent'' of height $1$ supported on $[\frac1{n+1},\frac1n]$. Then $\|f_n\|_\infty=1$, and $\|f_n-f_m\|_\infty=1$ for $n\ne m$ since the supports do not overlap:
no subsequence is Cauchy. The ball is complete but not totally bounded. This is no accident: F.~Riesz proved in 1918 that the closed unit
ball of a normed space is compact only in finite dimension (Tome~7). Arzel\`a--Ascoli (Section~\ref{sec:ascoli}) will say which sets of continuous
functions are compact.

\begin{theorem}[Heine]\label{thm:heine}
A continuous map $f:X\to Y$ from a compact metric space to a metric space is uniformly continuous: for every $\varepsilon>0$ there is $\delta>0$ with
$d(f(x),f(y))<\varepsilon$ whenever $d(x,y)<\delta$.
\end{theorem}
\begin{proof}
Otherwise, for some $\varepsilon>0$ and every $n$ there are $x_n,y_n$ with $d(x_n,y_n)<\frac1n$ and $d(f(x_n),f(y_n))\ge\varepsilon$. A subsequence $x_{n_k}$ converges to some $x$, and then
$y_{n_k}\to x$ too. By continuity at $x$, $f(x_{n_k})$ and $f(y_{n_k})$ both tend to $f(x)$, so $d(f(x_{n_k}),f(y_{n_k}))\to0$: a contradiction.
\end{proof}

$x\mapsto x^2$ on $\R$ and $x\mapsto\frac1x$ on $]0,1]$ are continuous and not uniformly continuous: compactness is needed. Heine's theorem is the key to
the existence of the integral of a continuous function (Chapter~7): on a fine enough subdivision, $f$ varies by less than $\varepsilon$ on every piece.

\section{Rolle and the mean value theorem}

Tome~0 admitted the mean value theorem; it rests on extreme values. A function $f$ defined near $c$ is \emph{differentiable} at $c$ if
$\frac{f(x)-f(c)}{x-c}$ has a limit $f'(c)$ as $x\to c$, $x\ne c$ (as throughout Tome~0).

\begin{theorem}\label{thm:mvt}
\begin{enumerate}[label=(\alph*)]
\item (Fermat) If $f:\,]a,b[\,\to\R$ has a maximum or a minimum at $c$ and is differentiable at $c$, then $f'(c)=0$.
\item (Rolle) If $f:[a,b]\to\R$ is continuous, differentiable on $]a,b[$, and $f(a)=f(b)$, then $f'(c)=0$ for some $c\in\,]a,b[$.
\item (Mean value theorem) If $f:[a,b]\to\R$ is continuous and differentiable on $]a,b[$, then $f(b)-f(a)=f'(c)(b-a)$ for some $c\in\,]a,b[$.
\item Consequently, on an interval $I$: if $f'\ge0$ on the interior, $f$ is non-decreasing; if $f'=0$, $f$ is constant; if $|f'|\le M$, $f$ is $M$-Lipschitz.
\end{enumerate}
\end{theorem}
\begin{proof}
(a) At a maximum, $\frac{f(x)-f(c)}{x-c}$ is $\le0$ for $x>c$ and $\ge0$ for $x<c$; the limit $f'(c)$ is therefore both $\le0$ and $\ge0$. For a minimum, apply this to $-f$.

(b) By Corollary~\ref{cor:extreme}, $f$ attains a maximum and a minimum on $[a,b]$. If both are attained at endpoints, they equal $f(a)=f(b)$, so $f$ is constant and
$f'=0$ everywhere. Otherwise one is attained at some $c\in\,]a,b[$, and $f'(c)=0$ by (a).

(c) Apply (b) to $g(x)=f(x)-f(a)-\frac{f(b)-f(a)}{b-a}(x-a)$, which has $g(a)=g(b)=0$ and $g'(x)=f'(x)-\frac{f(b)-f(a)}{b-a}$.

(d) For $x<y$ in $I$, (c) on $[x,y]$ gives $f(y)-f(x)=f'(c)(y-x)$ with $c$ interior to $I$. Each claim follows.
\end{proof}

With Theorem~\ref{thm:mvt}, every result of Tome~0 that relied on the mean value theorem (Taylor's formula with remainder, the
monotonicity of $\exp$, l'H\^opital's rule, the convexity inequalities) now rests on proved ground. The step in
Chapter~2 that used $|\cos x-\cos y|\le\sin1\,|x-y|$ is (d) with $M=\sin1$.

\begin{beauty}
Compactness turns ``for each point'' into ``uniformly''. Continuity at each point becomes uniform continuity (Heine); a positive
continuous function becomes bounded below by a positive constant (the proof of Corollary~\ref{cor:norms}); local Lipschitz bounds become global
ones; an open cover becomes finite. Almost every compactness argument is this one move: get something locally, cover,
extract a finite subcover, and take the worst of finitely many constants.
\end{beauty}

\section{Tychonoff's theorem and the Cantor space}

Proposition~\ref{prop:finiteproduct} handles finite products. For infinite products, the right tool is the finite intersection property and
the axiom of choice.

How can a product of uncountably many spaces, say $[0,1]^{\R}$, the set of \emph{all} functions from $\R$ to $[0,1]$, be compact?
Sequences are hopeless there, and no metric describes the topology. Andrey Tychonoff proved it for products of segments in
1930, at the age of twenty-four; the general theorem (1935) needs the axiom of choice, and Kelley's theorem shows that there is
no way around it. The proof below goes through Zorn's lemma: it builds a maximal family of closed sets that look as if they
should have a common point, and shows that they do.

\begin{theorem}[Tychonoff, 1930]\label{thm:tychonoff}
Any product $\prod_{i\in I}X_i$ of compact spaces is compact.
\end{theorem}
\begin{proof}
Let $X=\prod_iX_i$ and $p_i:X\to X_i$ the projections. By the reformulation at the start of the chapter, it is enough to show: if $\mathcal A$ is a family of subsets
of $X$ with the finite intersection property, there is a point $x$ in the closure of every $A\in\mathcal A$. (Apply this to a family of closed sets.)

\emph{A maximal family.} Order by inclusion the families $\mathcal D\supseteq\mathcal A$ of subsets of $X$ with the finite intersection property. The union of a chain of such
families has it too, since finitely many of its members lie in one family of the chain. By Zorn's lemma (Theorem~\ref{thm:zorn}; here the axiom of
choice is used), there is a maximal such family $\mathcal D$. Maximality gives two properties:
\begin{enumerate}[label=(\arabic*)]
\item a finite intersection of members of $\mathcal D$ is in $\mathcal D$ (adding it keeps the finite intersection property);
\item a set $B$ that meets every member of $\mathcal D$ is in $\mathcal D$ (by (1), $B$ meets every finite intersection of members, so adding $B$ keeps the property).
\end{enumerate}

\emph{The point.} For each $i$, the sets $p_i(D)$, $D\in\mathcal D$, have the finite intersection property in $X_i$, so by compactness of $X_i$ their closures have a common
point; choose one, $x_i$ (again a use of choice), and let $x=(x_i)_{i\in I}$.

\emph{It is in every closure.} Let $D\in\mathcal D$ and let $W=\bigcap_{i\in F}p_i^{-1}(U_i)$ be a basic open neighbourhood of $x$, with $F$ finite and $U_i$ open in $X_i$, $x_i\in U_i$.
For $i\in F$ and any $D'\in\mathcal D$, $U_i$ meets $p_i(D')$ since $x_i\in\overline{p_i(D')}$; so $p_i^{-1}(U_i)$ meets $D'$. By (2), $p_i^{-1}(U_i)\in\mathcal D$, and by (1), $W\in\mathcal D$. So $W$
meets $D$, by the finite intersection property. Every basic neighbourhood of $x$ meets $D$: $x\in\bar D$. Since $\mathcal A\subseteq\mathcal D$, we are done.
\end{proof}

\begin{remark}
Choice cannot be avoided: Kelley (1950) showed that Tychonoff's theorem implies the axiom of choice. For \emph{countable} products of compact
\emph{metric} spaces, it can: such a product is metrizable (Proposition~\ref{prop:product}(c)), and a diagonal extraction, as in the proof of
Theorem~\ref{thm:compactmetric}, shows it is sequentially compact (Exercise~\ref{ex:hilbertcube}).
\end{remark}

The most important compact product for this series is the space of infinite coin tosses.

\begin{definition}
The \emph{Cantor space} is $\{0,1\}^\N$ (here $\N=\{1,2,\dots\}$ indexes the tosses) with the product of discrete topologies. The \emph{Cantor set} is
\[C=\Bigl\{\sum_{k\ge1}\frac{2a_k}{3^k}:\ a\in\{0,1\}^\N\Bigr\}\subseteq[0,1],\]
the numbers with a ternary expansion using only the digits $0$ and $2$.
\end{definition}

A basic open set of $\{0,1\}^\N$ is a \emph{cylinder}: the set $[w]$ of sequences that begin with a given finite word $w$. The cylinders form the binary tree of
Figure~\ref{fig:cantortree}, and a point of $\{0,1\}^\N$ is an infinite branch.

\begin{figure}[ht]
\centering
\fig{ch04-cantortree}
\caption{The Cantor space as a binary tree. A word $w$ of length $n$ corresponds to a cylinder $[w]$ and to one of the $2^n$ intervals of length $3^{-n}$ that make up the
$n$-th stage of the Cantor set. An infinite branch ($0,1,1,\dots$ in orange) is a point; its intervals shrink to one point of $C$.}
\label{fig:cantortree}
\end{figure}

\begin{proposition}\label{prop:cantor}
\begin{enumerate}[label=(\alph*)]
\item $\{0,1\}^\N$ is compact, and metrizable by $\delta(a,b)=2^{-n}$, where $n$ is the first index with $a_n\ne b_n$ ($\delta(a,a)=0$).
\item $h(a)=\sum_k2a_k3^{-k}$ is a homeomorphism from $\{0,1\}^\N$ onto $C$. So $C$ is compact. Moreover $C=\bigcap_nC_n$, where $C_n$ is the union of the $2^n$ intervals
$I_w=[h(w00\ldots),h(w11\ldots)]$, of length $3^{-n}$, over the words $w$ of length $n$.
\item $g(a)=\sum_ka_k2^{-k}$ is a continuous surjection from $\{0,1\}^\N$ onto $[0,1]$. Hence $g\circ h^{-1}$ is a continuous surjection from $C$ onto $[0,1]$.
\end{enumerate}
\end{proposition}
\begin{proof}
(a) Compactness is Tychonoff's theorem ($\{0,1\}$ is finite, hence compact). The open $\delta$-ball of radius $2^{-n}$ around $a$ is the cylinder of the first $n$ terms of $a$:
so the $\delta$-open sets are the unions of cylinders, which are the open sets of the product topology. ($\delta$ is an ultrametric: if $a$ and $b$ agree up to
$n$, and $b$ and $c$ do, so do $a$ and $c$.)

(b) If $a$ and $b$ agree up to index $n$, $|h(a)-h(b)|\le\sum_{k>n}2\cdot3^{-k}=3^{-n}$: $h$ is continuous. If $n$ is the first index with $a_n\ne b_n$, the same estimate as in
Theorem~\ref{thm:continuum} gives $|h(a)-h(b)|\ge2\cdot3^{-n}-3^{-n}>0$: $h$ is injective. By Proposition~\ref{prop:compactbasic}(c) and (d), $h$ is a homeomorphism onto its
compact image $C$. If $a$ begins with $w$, then $h(w00\ldots)\le h(a)\le h(w11\ldots)$, and the difference of the ends is $\sum_{k>n}2\cdot3^{-k}=3^{-n}$: so $C\subseteq C_n$. The two
intervals $I_{w0}$ and $I_{w1}$ are the left and right thirds of $I_w$, so the intervals of distinct words of the same length are disjoint. If $x\in\bigcap_nC_n$, then
for each $n$, $x$ lies in exactly one $I_{w_n}$ with $|w_n|=n$, and $I_{w_{n+1}}\subseteq I_{w_n}$ forces $w_{n+1}$ to extend $w_n$. The words define a sequence $a$, and
$|x-h(a)|\le3^{-n}$ for every $n$ since both lie in $I_{w_n}$: $x=h(a)\in C$.

(c) $g$ is continuous by the same estimate with base $2$. Every $x\in[0,1]$ has a binary expansion: let $a_k$ be the $k$-th digit of the greedy expansion
($a_k=1$ iff $\sum_{j<k}a_j2^{-j}+2^{-k}\le x$), so that $0\le x-\sum_{j\le k}a_j2^{-j}\le2^{-k}$, and $g(a)=x$.
\end{proof}

The Cantor set is compact, uncountable (it is in bijection with $\{0,1\}^\N$), and has length zero: $C_n$ has total length $(\frac23)^n$. Chapter~6 makes ``length
zero'' precise. The surjection $g\circ h^{-1}$ extends to the Cantor staircase of Chapter~2 (Exercise~\ref{ex:devil}) and leads to a continuous surjection
$[0,1]\to[0,1]^2$, the curve that Cantor's bijection of Chapter~1 could not be (Exercise~\ref{ex:lebesguecurve}).

\section{Arzel\`a--Ascoli}\label{sec:ascoli}

In this section $K$ is a compact metric space and $C(K)$ the space of continuous real (or complex) functions on $K$ with the sup norm, a Banach space
(Proposition~\ref{prop:complete}(c); continuous functions on $K$ are bounded by Corollary~\ref{cor:extreme}). A subset of a metric space is \emph{relatively compact}
if its closure is compact.

\begin{definition}
A family $\mathcal F\subseteq C(K)$ is \emph{equicontinuous} at $x\in K$ if for every $\varepsilon>0$ there is $\delta>0$ such that $|f(y)-f(x)|<\varepsilon$ for all $f\in\mathcal F$ and all $y$ with $d(x,y)<\delta$: the
same $\delta$ serves every member of the family. It is \emph{equicontinuous} if it is equicontinuous at every point, and \emph{pointwise bounded} if
$\sup_{f\in\mathcal F}|f(x)|<\infty$ for each $x$.
\end{definition}

\begin{figure}[ht]
\centering
\fig{ch04-equicont}
\caption{Left: the functions $x\mapsto x^n$ on $[0,1]$ are not equicontinuous at $1$: each is continuous, but they need ever smaller $\delta$ there, and no
subsequence converges uniformly. Right: the functions $x\mapsto\frac12\sin(x+c)+\frac c4$, $0\le c\le2$, are all $\frac12$-Lipschitz, hence equicontinuous, and bounded: Arzel\`a--Ascoli
says the family is relatively compact.}
\label{fig:equicont}
\end{figure}

\begin{theorem}[Arzel\`a--Ascoli]\label{thm:ascoli}
A family $\mathcal F\subseteq C(K)$ is relatively compact iff it is equicontinuous and pointwise bounded.
\end{theorem}
\begin{proof}
($\Leftarrow$) Since $C(K)$ is complete, the closure $\bar{\mathcal F}$ is complete (Proposition~\ref{prop:complete}(a)); by Theorem~\ref{thm:compactmetric} it is enough to show that
$\mathcal F$ is totally bounded (then so is $\bar{\mathcal F}$: a $\varepsilon$-net of $\mathcal F$ is a $2\varepsilon$-net of $\bar{\mathcal F}$).

Let $\varepsilon>0$. For each $x\in K$ choose $\delta_x$ as in the definition of equicontinuity. Finitely many balls $B(x_j,\delta_{x_j})$, $1\le j\le m$, cover $K$. The vectors
\[v(f)=\bigl(f(x_1),\dots,f(x_m)\bigr),\qquad f\in\mathcal F,\]
lie in a bounded subset of $\R^m$ (or $\C^m$) with the norm $\|\cdot\|_\infty$, which is contained in a compact cube, hence is totally bounded: there are
$f_1,\dots,f_r\in\mathcal F$ such that every $f\in\mathcal F$ has some $l$ with $\|v(f)-v(f_l)\|_\infty<\varepsilon$. (A subset $E$ of a totally bounded space is totally bounded: if
$\varepsilon$-balls around $c_1,\dots,c_p$ cover the space, pick a point of $E$ in each ball that meets $E$; the $2\varepsilon$-balls around these points cover $E$.) Now for $y\in K$ pick $j$
with $d(y,x_j)<\delta_{x_j}$; then
\[\begin{aligned}|f(y)-f_l(y)|&\le|f(y)-f(x_j)|+|f(x_j)-f_l(x_j)|\eqbr+|f_l(x_j)-f_l(y)|<3\varepsilon.\end{aligned}\]
So $\|f-f_l\|_\infty<3\varepsilon$: the $f_l$ form a $3\varepsilon$-net of $\mathcal F$.

($\Rightarrow$) If $\bar{\mathcal F}$ is compact it is bounded, so $\mathcal F$ is pointwise bounded. Let $x\in K$ and $\varepsilon>0$; take an $\varepsilon$-net $f_1,\dots,f_r$ of $\mathcal F$ and $\delta>0$ such
that $|f_l(y)-f_l(x)|<\varepsilon$ for $d(x,y)<\delta$ and all $l$ (finitely many conditions). For $f\in\mathcal F$, choose $l$ with $\|f-f_l\|_\infty<\varepsilon$; then $|f(y)-f(x)|<3\varepsilon$ when
$d(x,y)<\delta$.
\end{proof}

\begin{corollary}\label{cor:lipcompact}
Let $L,M\ge0$ and $x_0\in K$. The set of $L$-Lipschitz functions $f:K\to\R$ with $|f(x_0)|\le M$ is compact in $C(K)$. In particular every sequence of such
functions has a uniformly convergent subsequence.
\end{corollary}
\begin{proof}
The family is equicontinuous (take $\delta=\varepsilon/L$, or any $\delta$ if $L=0$) and bounded by $M+L\operatorname{diam}K$. It is closed, since a uniform (even pointwise) limit of $L$-Lipschitz
functions is $L$-Lipschitz, with $|f(x_0)|\le M$. Theorem~\ref{thm:ascoli} applies.
\end{proof}

This is the classical use of Arzel\`a--Ascoli: to produce a solution of a problem as the limit of approximate solutions, without uniqueness and
without a formula. Peano's existence theorem for differential equations is the model (Exercise~\ref{ex:peano}); in probability, the same idea
constructs Brownian motion as a limit of random walks (Tome~11).

\section{Stone--Weierstrass}

In 1885 Karl Weierstrass, then seventy years old, proved that every continuous function on $[a,b]$ is a uniform limit of polynomials. It
was a surprise: continuous functions can be wild (Chapter~5 shows that most of them have no derivative anywhere), and polynomials are the
tamest functions there are. Yet the tame ones approximate the wild ones as closely as we wish. In 1937, and in a simpler form in 1948,
Marshall Stone found the reason, and it has nothing to do with polynomials in particular: they form an algebra that separates points. In this section $K$ is a compact space and $C(K,\R)$, $C(K,\C)$ the continuous functions
with the sup norm. A uniform limit of continuous functions on $K$ is continuous, by the proof of Proposition~\ref{prop:complete}(c), which never used a
metric on $K$.

\begin{definition}
A set $\mathcal A$ of functions on $K$ \emph{separates points} if for all $x\ne y$ in $K$ some $f\in\mathcal A$ has $f(x)\ne f(y)$. A \emph{subalgebra} of $C(K,\R)$ is a linear subspace
stable under products.
\end{definition}

The proof needs a single concrete approximation: of $|t|$ by polynomials.

\begin{lemma}\label{lem:abs}
Let $p_0=0$ and $p_{n+1}(t)=p_n(t)+\frac12\bigl(t^2-p_n(t)^2\bigr)$. Each $p_n$ is a polynomial, and $0\le|t|-p_n(t)\le\frac2{en}$ for $t\in[-1,1]$ and $n\ge1$.
\end{lemma}
\begin{proof}
Write $s=|t|\in[0,1]$. Then $s-p_{n+1}=(s-p_n)\bigl(1-\frac{s+p_n}2\bigr)$, as one checks by expanding ($t^2=s^2$). If $0\le p_n\le s$, the second factor is in
$[1-s,1]\subseteq[0,1]$, so $0\le s-p_{n+1}\le s-p_n$, and $p_{n+1}\ge p_n\ge0$. By induction, $0\le s-p_n$ and, since $p_n\ge0$, $s-p_{n+1}\le(s-p_n)(1-\frac s2)$, so
\[0\le s-p_n\le s\Bigl(1-\frac s2\Bigr)^n\le se^{-ns/2}\le\frac2{en},\]
using $1-u\le e^{-u}$ and the maximum of $s\mapsto se^{-ns/2}$, attained at $s=\frac2n$.
\end{proof}

\begin{theorem}[Stone--Weierstrass]\label{thm:sw}
\begin{enumerate}[label=(\alph*)]
\item (Real form) A subalgebra $\mathcal A\subseteq C(K,\R)$ that contains the constants and separates points is dense: every $f\in C(K,\R)$ is a uniform limit of
elements of $\mathcal A$.
\item (Complex form) A subalgebra $\mathcal A\subseteq C(K,\C)$ that contains the constants, separates points and is stable under complex conjugation is
dense in $C(K,\C)$.
\end{enumerate}
\end{theorem}
\begin{proof}
(a) \emph{Step 1: the closure $\bar{\mathcal A}$ is a subalgebra, stable under $|\cdot|$, $\max$ and $\min$.} If $f_n\to f$ and $g_n\to g$ uniformly, then $f_n+g_n\to f+g$, $cf_n\to cf$, and
$\|f_ng_n-fg\|\le\|f_n\|\|g_n-g\|+\|g\|\|f_n-f\|\to0$. If $f\in\bar{\mathcal A}$ and $M=\|f\|_\infty>0$, then $Mp_n(f/M)\in\bar{\mathcal A}$ (a polynomial in $f$) and
$\bigl|Mp_n(f/M)-|f|\bigr|\le\frac{2M}{en}$ by Lemma~\ref{lem:abs}; so $|f|\in\bar{\mathcal A}$. Then $\max(f,g)=\frac12(f+g+|f-g|)$ and $\min(f,g)=\frac12(f+g-|f-g|)$ are in $\bar{\mathcal A}$.

\emph{Step 2: two-point interpolation.} For $x\ne y$ and real $a,b$, some $h\in\mathcal A$ has $h(x)=a$, $h(y)=b$: take $u\in\mathcal A$ with $u(x)\ne u(y)$ and
\[h=a+(b-a)\frac{u-u(x)}{u(y)-u(x)}.\]

\emph{Step 3: from below.} Let $f\in C(K,\R)$ and $\varepsilon>0$, and fix $x\in K$. For each $y$, Step~2 gives $h_y\in\mathcal A$ with $h_y(x)=f(x)$ and $h_y(y)=f(y)$. The set $V_y=\{h_y<f+\varepsilon\}$
is open and contains $y$ (and $x$). Finitely many $V_{y_1},\dots,V_{y_m}$ cover $K$; then $g_x=\min(h_{y_1},\dots,h_{y_m})\in\bar{\mathcal A}$ satisfies $g_x<f+\varepsilon$ everywhere and $g_x(x)=f(x)$.

\emph{Step 4: from above.} The set $W_x=\{g_x>f-\varepsilon\}$ is open and contains $x$. Finitely many $W_{x_1},\dots,W_{x_r}$ cover $K$; then $g=\max(g_{x_1},\dots,g_{x_r})\in\bar{\mathcal A}$
satisfies $f-\varepsilon<g<f+\varepsilon$ everywhere. So $\|f-g\|_\infty\le\varepsilon$ (in fact $<\varepsilon$, by compactness), and $f\in\bar{\mathcal A}$.

(b) Let $\mathcal A_\R$ be the real-valued functions in $\mathcal A$. If $f\in\mathcal A$, then $\operatorname{Re}f=\frac12(f+\bar f)$ and $\operatorname{Im}f=\frac1{2i}(f-\bar f)$ are in $\mathcal A_\R$. So $\mathcal A_\R$ is a real
subalgebra containing the real constants, and it separates points: if $f(x)\ne f(y)$, then $\operatorname{Re}f$ or $\operatorname{Im}f$ separates $x$ and $y$. By (a), $\mathcal A_\R$ is dense in
$C(K,\R)$, and every $f=u+iv\in C(K,\C)$ is approximated by $u_n+iv_n$ with $u_n,v_n\in\mathcal A_\R$.
\end{proof}

\begin{think}
Why does the complex form need stability under conjugation? Try to approximate $z\mapsto\bar z$ on the unit circle by polynomials in $z$.
\end{think}

It cannot be done. For a polynomial $p(z)=\sum_{k=0}^dc_kz^k$,
\[\int_0^{2\pi}p(e^{it})\,e^{it}\,dt=\sum_kc_k\int_0^{2\pi}e^{i(k+1)t}\,dt=0,\]
since $e^{i(k+1)t}$ has primitive $\frac{e^{i(k+1)t}}{i(k+1)}$, which is $2\pi$-periodic. If $p_n(z)\to\bar z$ uniformly on the circle, the integrals would tend to
$\int_0^{2\pi}e^{-it}e^{it}\,dt=2\pi$. Polynomials in $z$ separate the points of the circle and contain the constants; by the maximum principle (Tome~4), their uniform limits on the circle
extend to continuous functions on the closed disc that are holomorphic inside, a much smaller class. Conjugation is what lets the complex
algebra see $\operatorname{Re}$ and $\operatorname{Im}$.

\begin{corollary}\label{cor:weierstrass}
\begin{enumerate}[label=(\alph*)]
\item (Weierstrass, 1885) For compact $K\subseteq\R^n$, every continuous $f:K\to\R$ is a uniform limit of polynomials in $x_1,\dots,x_n$.
\item (Trigonometric form) Every continuous $2\pi$-periodic $f:\R\to\C$ is a uniform limit of trigonometric polynomials $\sum_{|k|\le N}c_ke^{ikt}$.
\item $C([a,b])$ is separable: the polynomials with rational coefficients are dense.
\end{enumerate}
\end{corollary}
\begin{proof}
(a) The polynomials form a subalgebra containing the constants, and the coordinate functions separate points.

(b) By Theorem~\ref{thm:exp}(e), $t\mapsto e^{it}$ maps $[0,2\pi]$ continuously onto the circle $S^1$, identifying only $0$ and $2\pi$. A $2\pi$-periodic $f$ therefore defines
$F:S^1\to\C$ with $f(t)=F(e^{it})$, and $F$ is continuous: if $F$ were discontinuous at some point $e^{is}$, a sequence $e^{it_n}\to e^{is}$ with $|F(e^{it_n})-F(e^{is})|\ge\varepsilon$ could
be taken with $t_n\in[0,2\pi]$; a subsequence $t_n\to t$ with $e^{it}=e^{is}$, and $f(t_n)\to f(t)=F(e^{is})$ by continuity of $f$, a contradiction. The polynomials in $z$ and $\bar z$
on $S^1$ form a subalgebra of $C(S^1,\C)$ with the constants, stable under conjugation, and $z$ separates points. On $S^1$, $z^k\bar z^l=e^{i(k-l)t}$. Theorem~\ref{thm:sw}(b)
applies.

(c) Approximate $f$ within $\varepsilon$ by a polynomial $p$, then the coefficients of $p$ by rationals: on $[a,b]$, changing a coefficient of $x^k$ by $\eta$ changes $p$ by at most
$\eta\max(|a|,|b|)^k$.
\end{proof}

Weierstrass's own proof convolved $f$ with a Gaussian, an approximate identity (Chapter~9). Bernstein's 1912 proof is probabilistic and
constructive: it approximates $f$ by $\mathbb E\bigl[f(\frac{S_n}n)\bigr]$ where $S_n$ counts heads in $n$ tosses of a coin with bias $x$ (Exercise~\ref{ex:bernstein}). Part~(b) is the
starting point of Fourier analysis in $L^2$ (Tome~7); Fej\'er's theorem gives explicit trigonometric approximants (Chapter~9).

\begin{remark}
Stone--Weierstrass explains why one-hidden-layer neural networks are \emph{universal approximators}. For instance, the functions
$x\mapsto\sum_jc_j\exp(w_j\cdot x)$ on a compact $K\subseteq\R^n$ form an algebra ($e^{u}e^{v}=e^{u+v}$) with the constants ($w=0$) that separates points, so they are dense in $C(K,\R)$.
Cybenko (1989) and Hornik (1991) proved the same for sums $\sum_jc_j\sigma(w_j\cdot x+b_j)$ with any sigmoid $\sigma$; for the ReLU $\sigma(t)=\max(t,0)$ in one variable the sums
are exactly the continuous piecewise linear functions, dense by Heine's theorem. A numerical calculation trains such networks. Density says nothing about
how many neurons are needed, nor whether training finds them.
\end{remark}

\section{Exercises}

\begin{exercise}[$\star$]\label{ex:dini}
\emph{Dini's theorem.} Let $K$ be compact and $f_n,f\in C(K,\R)$ with $f_1\ge f_2\ge\cdots$ and $f_n(x)\to f(x)$ for every $x$. Show that $f_n\to f$ uniformly. Show by an example that each
hypothesis (compactness, monotonicity, continuity of $f$) is needed.
\end{exercise}

\begin{exercise}[$\star$]\label{ex:moments}
Let $f:[0,1]\to\R$ be continuous with $\int_0^1f(x)x^n\,dx=0$ for every $n\ge0$. Show that $f=0$. (The moments of a continuous function determine it.)
\end{exercise}

\begin{exercise}[$\star\star$]\label{ex:isometry}
Let $K$ be a compact metric space and $f:K\to K$ with $d(f(x),f(y))=d(x,y)$ for all $x,y$. Show that $f$ is surjective. Find a complete, bounded metric space where this fails.
\end{exercise}

\begin{exercise}[$\star\star$]\label{ex:edelstein}
\emph{Edelstein's fixed point theorem.} Let $K$ be a compact metric space and $f:K\to K$ with $d(f(x),f(y))<d(x,y)$ for $x\ne y$. Show that $f$ has a unique fixed point $p$
and that $f^n(x)\to p$ for every $x$. Check that $\sin$ on $[0,1]$ satisfies the hypothesis without being a contraction in the sense of Banach's theorem.
\end{exercise}

\begin{exercise}[$\star\star$]\label{ex:hilbertcube}
\emph{The Hilbert cube.} Let $Q=[0,1]^\N$ with the metric $d(x,y)=\sum_k2^{-k}|x_k-y_k|$. (a) Show, by a diagonal extraction and without Tychonoff's theorem, that $Q$ is
compact. (b) Let $(X,d)$ be a compact metric space of diameter at most $1$ and $(a_k)$ a dense sequence. Show that $\iota(x)=(d(x,a_k))_{k\ge1}$ is a homeomorphism of $X$ onto a
compact subset of $Q$. Every compact metric space lives in the Hilbert cube.
\end{exercise}

\begin{exercise}[$\star\star$]\label{ex:bernstein}
\emph{Bernstein's proof of Weierstrass.} For $f\in C([0,1])$ let $B_nf(x)=\sum_{k=0}^nf(\frac kn)\binom nkx^k(1-x)^{n-k}$. (a) Show that $\sum_k\binom nkx^k(1-x)^{n-k}\bigl(\frac kn-x\bigr)^2=\frac{x(1-x)}n$.
(b) Deduce that $B_nf\to f$ uniformly on $[0,1]$. (c) Interpret $B_nf(x)$ with $n$ tosses of a coin with probability $x$ of heads.
\end{exercise}

\begin{exercise}[$\star\star$]\label{ex:lebesguecurve}
\emph{Lebesgue's space-filling curve.} (a) Using Proposition~\ref{prop:cantor}, build a continuous surjection $\psi:C\to[0,1]^2$. (Split a sequence into its odd and even terms.)
(b) Extend $\psi$ to $\Psi:[0,1]\to[0,1]^2$, affine on each open interval of $[0,1]\setminus C$, and show that $\Psi$ is continuous. So $[0,1]$ maps continuously \emph{onto} the square, although
not bijectively (Exercise~\ref{ex:notR2}).
\end{exercise}

\begin{exercise}[$\star\star\star$]\label{ex:peano}
\emph{Peano's existence theorem.} Let $f:[0,1]\times\R\to\R$ be continuous and bounded by $M$, and $y_0\in\R$. (a) For $n\ge1$ let $y_n$ be the \emph{Euler polygon}: $y_n(0)=y_0$ and
on each $[\frac kn,\frac{k+1}n]$, $y_n$ is affine with slope $f\bigl(\frac kn,y_n(\frac kn)\bigr)$. Show that a subsequence converges uniformly to some $y$. (b) Show that
$y(t)=y_0+\int_0^tf(s,y(s))\,ds$, so that $y$ is a $C^1$ solution of $y'=f(t,y)$, $y(0)=y_0$. (c) For $f(t,y)=\sqrt{\min(|y|,1)}$ and $y_0=0$, find two solutions. Existence does not need
the Lipschitz condition of Picard--Lindel\"of (Exercise~\ref{ex:picard}); uniqueness does.
\end{exercise}

\section{Hints}

\begin{hint}{ex:dini}
Look at the closed sets $F_n=\{f_n-f\ge\varepsilon\}$. For the examples, think of $x^n$ on $[0,1[$ and on $[0,1]$.
\end{hint}
\begin{hint}{ex:moments}
$\int_0^1fp=0$ for every polynomial $p$. Approximate $f$ itself by polynomials and look at $\int_0^1f^2$.
\end{hint}
\begin{hint}{ex:isometry}
If $y\notin f(K)$, its distance $\varepsilon$ to the compact set $f(K)$ is positive. Look at the orbit $y,f(y),f^2(y),\dots$.
\end{hint}
\begin{hint}{ex:edelstein}
Minimize the continuous function $x\mapsto d(x,f(x))$. For convergence, $d(f^n(x),p)$ decreases; look at a limit point of the orbit.
\end{hint}
\begin{hint}{ex:hilbertcube}
(a) Extract a subsequence converging in the first coordinate, then a further one in the second, and so on, and take the diagonal. (b) Each
coordinate of $\iota$ is $1$-Lipschitz; to separate $x\ne y$, take $a_k$ very close to $x$.
\end{hint}
\begin{hint}{ex:bernstein}
(a) Differentiate $(x+y)^n=\sum\binom nkx^ky^{n-k}$ twice, or compute $\sum k\binom nk\ldots$ and $\sum k(k-1)\binom nk\ldots$. (b) Split the sum into $|\frac kn-x|<\delta$ and the rest, and use Heine.
\end{hint}
\begin{hint}{ex:lebesguecurve}
(a) $\psi=(g\circ\mathrm{odd},g\circ\mathrm{even})\circ h^{-1}$. (b) Only finitely many gaps are longer than any given $\eta>0$. For $x_n\to x\in C$ with $x_n$ in gaps, the endpoints
of those gaps tend to $x$ too, unless one gap is visited infinitely often.
\end{hint}
\begin{hint}{ex:peano}
(a) Each $y_n$ is $M$-Lipschitz with $y_n(0)=y_0$: use Corollary~\ref{cor:lipcompact}. (b) Write $y_n(t)=y_0+\int_0^tf\bigl(\frac{\lfloor ns\rfloor}n,y_n(\frac{\lfloor ns\rfloor}n)\bigr)\,ds$ and use the uniform continuity of $f$ on a
compact rectangle. (c) Try $y=0$ and $y=\frac{t^2}4$.
\end{hint}

\section{Solutions}

\begin{solution}{ex:dini}
\idea Pointwise convergence says each point eventually behaves; monotonicity says that once a point behaves, it stays good; compactness turns
``each point eventually'' into ``all points at once'' through the finite intersection property.
\writeup Let $\varepsilon>0$ and $F_n=\{x:f_n(x)-f(x)\ge\varepsilon\}$, closed since $f_n-f$ is continuous. Since $f_n\ge f_{n+1}\ge f$, the $F_n$ decrease, and their intersection is empty since
$f_n(x)\to f(x)$. By compactness some finite intersection is empty, that is, some $F_N=\varnothing$, and then $F_n=\varnothing$ for $n\ge N$: $0\le f_n-f<\varepsilon$ everywhere. Examples: $x^n\to0$
on $[0,1[$ decreasingly, not uniformly (not compact); on $[0,1]$, $x^n$ decreases to the discontinuous indicator of $\{1\}$; and the tents $f_n$ of height $1$ on
$[\frac1{n+1},\frac1n]$ tend to $0$ pointwise on $[0,1]$ without being monotone, and not uniformly.
\end{solution}

\begin{solution}{ex:moments}
\idea If $f$ is orthogonal to all polynomials, it is orthogonal to everything polynomials approximate, including $f$ itself.
\writeup By linearity $\int_0^1fp=0$ for every polynomial $p$. Choose polynomials $p_n\to f$ uniformly on $[0,1]$ (Corollary~\ref{cor:weierstrass}). Then
\[\Bigl|\int_0^1f^2-\int_0^1fp_n\Bigr|\le\|f\|_\infty\|f-p_n\|_\infty\to0,\]
so $\int_0^1f^2=0$. A continuous non-negative function with integral $0$ vanishes (if $f^2(x_0)>0$, it is $>\frac12f^2(x_0)$ on an interval around $x_0$). So $f=0$.
\end{solution}

\begin{solution}{ex:isometry}
\idea An isometry that misses a point pushes it along an orbit whose points stay at distance at least $\varepsilon$ from each other; a compact space
has no room for infinitely many such points.
\writeup $f$ is continuous, so $f(K)$ is compact, hence closed. If $y\notin f(K)$, then $\varepsilon=d(y,f(K))>0$. Let $y_n=f^n(y)$. For $m>n$, since $f^n$ is an isometry,
$d(y_n,y_m)=d(y,f^{m-n}(y))\ge\varepsilon$, as $f^{m-n}(y)\in f(K)$. No subsequence of $(y_n)$ is Cauchy, contradicting sequential compactness
(Theorem~\ref{thm:compactmetric}). Without compactness: on $\N$ with the discrete metric ($d=1$ between distinct points), complete and bounded,
$n\mapsto n+1$ is an isometry that misses $0$.
\end{solution}

\begin{solution}{ex:edelstein}
\idea A strict inequality between distances is not enough to iterate Banach's argument, since the ratio may tend to $1$. Compactness replaces
the missing uniform constant: a continuous function attains its minimum, and orbits have limit points.
\writeup $\phi(x)=d(x,f(x))$ is continuous ($f$ is $1$-Lipschitz), so it attains its minimum at some $p$. If $f(p)\ne p$, then $\phi(f(p))=d(f(p),f(f(p)))<d(p,f(p))=\phi(p)$, a
contradiction: $f(p)=p$. If $q$ is another fixed point, $d(p,q)=d(f(p),f(q))<d(p,q)$ is absurd. Let $x\in K$, $x_n=f^n(x)$ and $r_n=d(x_n,p)$. Since $d(f(x_n),f(p))\le d(x_n,p)$,
$(r_n)$ is non-increasing, with limit $r\ge0$. A subsequence $x_{n_k}\to q$, and $d(q,p)=r$; also $x_{n_k+1}=f(x_{n_k})\to f(q)$, so $d(f(q),p)=\lim r_{n_k+1}=r$. If $r>0$ then $q\ne p$ and
$d(f(q),p)=d(f(q),f(p))<d(q,p)=r$, a contradiction. So $r=0$. For $\sin$ on $[0,1]$: $|\sin x-\sin y|=|\cos c|\,|x-y|<|x-y|$ for $x\ne y$ by the mean value theorem, since
$0<|\cos c|<1$ for $c\in\,]0,1]$; but $\frac{\sin x}x\to1$ as $x\to0$, so no constant $k<1$ works. The fixed point is $0$, approached very slowly: $x_n\sim\sqrt{3/n}$.
\end{solution}

\begin{solution}{ex:hilbertcube}
\idea Convergence in $Q$ is convergence of each coordinate, because the weights $2^{-k}$ make the tail negligible. So a diagonal subsequence that
converges coordinate by coordinate converges in $Q$. For (b), the distances to a dense sequence are coordinates that know where a point is.
\writeup (a) If $x^{(n)}_k\to x_k$ for every $k$, then $d(x^{(n)},x)\le\sum_{k\le K}|x^{(n)}_k-x_k|+2^{-K}$, which is $<2^{1-K}$ for $n$ large; so coordinatewise convergence implies
convergence in $Q$. Given a sequence $(x^{(n)})$ in $Q$, Theorem~\ref{thm:heineborel} gives an infinite set $N_1$ of indices along which the first coordinates converge, then
an infinite $N_2\subseteq N_1$ along which the second coordinates converge, and so on. Choosing $n_1<n_2<\cdots$ with $n_j\in N_j$, the sequence $(x^{(n_j)})_j$ lies in $N_k$ from rank $k$
on, so its $k$-th coordinates converge, for every $k$. It converges in $Q$, which is therefore sequentially compact, hence compact.
(b) $|d(x,a_k)-d(y,a_k)|\le d(x,y)$, so $d(\iota x,\iota y)\le d(x,y)$: $\iota$ is continuous, with values in $Q$ since $d\le1$. If $x\ne y$, let $r=d(x,y)$ and choose $a_k$ with $d(x,a_k)<\frac r3$;
then $d(y,a_k)>\frac{2r}3$, so $\iota x\ne\iota y$. A continuous injection of a compact space into a Hausdorff space is a homeomorphism onto its image
(Proposition~\ref{prop:compactbasic}(d)), which is compact. A dense sequence exists since compact metric spaces are separable.
\end{solution}

\begin{solution}{ex:bernstein}
\idea Under the binomial weights, $\frac kn$ concentrates around $x$ with variance $\frac{x(1-x)}n\le\frac1{4n}$: so the weighted average of the values $f(\frac kn)$ only sees values of $f$ near
$x$, where $f$ is close to $f(x)$ by uniform continuity.
\writeup (a) Let $w_k=\binom nkx^k(1-x)^{n-k}$. The binomial formula gives $\sum_kw_k=1$; from $k\binom nk=n\binom{n-1}{k-1}$ and $k(k-1)\binom nk=n(n-1)\binom{n-2}{k-2}$, $\sum_kkw_k=nx$ and
$\sum_kk(k-1)w_k=n(n-1)x^2$. So
\[\begin{aligned}\sum_kw_k\Bigl(\frac kn-x\Bigr)^2&=\frac{n(n-1)x^2+nx}{n^2}-2x^2+x^2\eqbr=\frac{x-x^2}n.\end{aligned}\]
(b) Let $\varepsilon>0$ and $\delta$ from Heine's theorem: $|f(s)-f(x)|<\varepsilon$ when $|s-x|<\delta$. Then
\[\begin{aligned}|B_nf(x)-f(x)|&\le\sum_kw_k\Bigl|f\Bigl(\frac kn\Bigr)-f(x)\Bigr|\\&\le\varepsilon+2\|f\|_\infty\sum_{|\frac kn-x|\ge\delta}w_k\eqbr\le\varepsilon+\frac{2\|f\|_\infty}{\delta^2}\cdot\frac1{4n},\end{aligned}\]
since on the second set $1\le(\frac kn-x)^2/\delta^2$. For $n\ge\|f\|_\infty/(2\delta^2\varepsilon)$, $|B_nf-f|\le2\varepsilon$ on $[0,1]$.
(c) If $S_n$ counts heads in $n$ independent tosses with probability $x$ of heads, $\mathbb P(S_n=k)=w_k$, so $B_nf(x)=\mathbb E\bigl[f(\frac{S_n}n)\bigr]$. Part (a) says $\operatorname{Var}(\frac{S_n}n)=\frac{x(1-x)}n$, and (b) is
Chebyshev's inequality: the law of large numbers, uniformly in $x$, proves Weierstrass's theorem.
\end{solution}

\begin{solution}{ex:lebesguecurve}
\idea A point of $C$ is a sequence of coin tosses; the odd tosses and the even tosses are two independent sequences, each of which names a point of
$[0,1]$. On the gaps of $C$ there is nothing to define, so join the ends by segments.
\writeup (a) The map $s:\{0,1\}^\N\to\{0,1\}^\N\times\{0,1\}^\N$, $a\mapsto\bigl((a_1,a_3,\ldots),(a_2,a_4,\ldots)\bigr)$, is a bijection, continuous since each output coordinate is an input coordinate.
Let $\psi=(g\times g)\circ s\circ h^{-1}$, with $g$, $h$ from Proposition~\ref{prop:cantor}: $h^{-1}$ is continuous by Proposition~\ref{prop:cantor}(b), $g\times g$ is continuous and onto
$[0,1]^2$, and $s$ is onto. So $\psi$ is a continuous surjection $C\to[0,1]^2$.

(b) $[0,1]\setminus C$ is a countable union of disjoint open intervals $]\alpha,\beta[$ (the ``gaps''), with $\alpha,\beta\in C$: by Proposition~\ref{prop:cantor}(b), $[0,1]\setminus C=\bigcup_n([0,1]\setminus C_n)$,
and each gap removed at stage $n$ has its ends in $C$. Define $\Psi=\psi$ on $C$ and $\Psi(x)=\psi(\alpha)+\frac{x-\alpha}{\beta-\alpha}(\psi(\beta)-\psi(\alpha))$ on a gap. $\Psi$ is continuous on each open gap.
Let $x\in C$ and $x_n\to x$; we show $\Psi(x_n)\to\Psi(x)$, and it is enough to show that every subsequence has a further subsequence along which this holds, and every subsequence has a further
subsequence of one of the following three kinds. If all $x_n\in C$, this is continuity of $\psi$.
If all $x_n$ lie in one gap $]\alpha,\beta[$, then $x\in\{\alpha,\beta\}$ (it is in $C$ and in $[\alpha,\beta]$), and the affine formula is continuous up to the ends. Otherwise, passing to
a subsequence, the $x_n$ lie in gaps $]\alpha_n,\beta_n[$ each visited finitely often. Only finitely many gaps have length $\ge\eta$ (stage $n$ removes gaps of length
$3^{-n}$), so $\beta_n-\alpha_n\to0$. As $x\notin\,]\alpha_n,\beta_n[$, $|\alpha_n-x|$ and $|\beta_n-x|$ are at most $|x_n-x|+(\beta_n-\alpha_n)\to0$. By continuity of $\psi$, $\psi(\alpha_n)$ and $\psi(\beta_n)$
tend to $\psi(x)$, and $\Psi(x_n)$, which lies on the segment between them, tends to $\psi(x)$ too. So $\Psi$ is continuous, and onto $[0,1]^2$ since $\psi$ already is.
\end{solution}

\begin{solution}{ex:peano}
\idea Euler's method produces approximate solutions that are all $M$-Lipschitz; Arzel\`a--Ascoli extracts a limit, and the uniform continuity
of $f$ on a compact set lets the limit pass inside the integral. Nothing forces different subsequences to have the same limit, and indeed
uniqueness can fail.
\writeup (a) Each $y_n$ is continuous, piecewise affine with slopes bounded by $M$, hence $M$-Lipschitz, and $y_n(0)=y_0$. By Corollary~\ref{cor:lipcompact} on $K=[0,1]$, a
subsequence $y_{n_j}$ converges uniformly to some $y$, itself $M$-Lipschitz.
(b) On $]\frac kn,\frac{k+1}n[$ the slope of $y_n$ is $f(\tau_n(s),y_n(\tau_n(s)))$ with $\tau_n(s)=\frac{\lfloor ns\rfloor}n$, so $y_n(t)=y_0+\int_0^tf\bigl(\tau_n(s),y_n(\tau_n(s))\bigr)\,ds$ (a step function,
integrated piece by piece). All values $y_n(s)$ and $y(s)$ lie in $[y_0-M,y_0+M]$, and $f$ is uniformly continuous on the compact $R=[0,1]\times[y_0-M,y_0+M]$ (Heine).
Since $|\tau_n(s)-s|\le\frac1n$ and
\[\begin{aligned}|y_n(\tau_n(s))-y(s)|&\le|y_n(\tau_n(s))-y_n(s)|+|y_n(s)-y(s)|\eqbr\le\tfrac Mn+\|y_n-y\|_\infty,\end{aligned}\]
along the subsequence the integrand tends to $f(s,y(s))$ uniformly in $s$. Passing to the limit, $y(t)=y_0+\int_0^tf(s,y(s))\,ds$. The integrand is continuous, so by the
fundamental theorem of calculus (Tome~0) $y$ is $C^1$ with $y'(t)=f(t,y(t))$ and $y(0)=y_0$.
(c) $y=0$ is a solution. For $y=\frac{t^2}4\le\frac14$ on $[0,1]$, $y'=\frac t2=\sqrt{y}=f(t,y)$. Both start at $0$. Here $f$ is not Lipschitz near $y=0$, since $\sqrt{|y|}/|y|\to\infty$.
\end{solution}

\section*{Sources and further reading}

\begin{itemize}
\item J.~Munkres, \emph{Topology}, 2nd ed., Prentice Hall, 2000, Chapters~3 and~5. Compactness, the tube lemma, and Tychonoff's
theorem by maximal families with the finite intersection property, as here.
\item W.~Rudin, \emph{Principles of Mathematical Analysis}, 3rd ed., McGraw--Hill, 1976, Chapters~2 and~7. Compact metric spaces,
equicontinuity, and Stone--Weierstrass; the iteration of Lemma~\ref{lem:abs} is his Exercise~7.23.
\item A.~Tychonoff, \emph{\"Uber die topologische Erweiterung von R\"aumen}, Mathematische Annalen 102 (1930), 544--561; J.~L.~Kelley,
\emph{The Tychonoff product theorem implies the axiom of choice}, Fundamenta Mathematicae 37 (1950), 75--76.
\item K.~Weierstrass, \emph{\"Uber die analytische Darstellbarkeit sogenannter willk\"urlicher Functionen einer reellen Ver\"anderlichen},
Sitzungsberichte der Akademie zu Berlin (1885), 633--639 and 789--805; M.~H.~Stone, \emph{The generalized Weierstrass approximation
theorem}, Mathematics Magazine 21 (1948), 167--184 and 237--254.
\item S.~Bernstein, \emph{D\'emonstration du th\'eor\`eme de Weierstrass fond\'ee sur le calcul des probabilit\'es}, Communications de la
Soci\'et\'e math\'ematique de Kharkow 13 (1912), 1--2.
\item H.~Sagan, \emph{Space-Filling Curves}, Springer, 1994. Lebesgue's curve, Peano's and Hilbert's.
\item G.~Cybenko, \emph{Approximation by superpositions of a sigmoidal function}, Mathematics of Control, Signals and Systems 2 (1989),
303--314; K.~Hornik, \emph{Approximation capabilities of multilayer feedforward networks}, Neural Networks 4 (1991), 251--257.
\item \emph{For the stories in the introduction.} M.~Raman-Sundstr\"om, ``A pedagogical history of compactness'', \emph{American Math. Monthly} 122 (2015), 619--635, follows
Heine, Borel, Fr\'echet and Alexandroff--Urysohn.
\end{itemize}
